\documentclass[11pt,a4paper]{article}

\usepackage[margin=2.5cm]{geometry}
\usepackage[T1]{fontenc}
\usepackage{lmodern}
\usepackage{parskip}
\usepackage{enumitem}
\usepackage{xurl}
\usepackage[hidelinks]{hyperref}

\begin{document}

\begin{center}
  {\Large\bfseries BAScii Generative AI -- Final Project Report}
\end{center}

\vspace{1em}

\begin{tabular}{@{}lp{11.5cm}@{}}
  \textbf{Project Title:}      & Yapperton \\
  \textbf{Group:}              & Aliens \\
  \textbf{Team Members:}       & Wilhelm Axel Durelius (6901009256) \\
                               & Shayan Gardan Gardawan (6901005756) \\
                               & Schahab Bashar-Mal (6901001156) \\
  \textbf{Project Repository:} & \url{https://github.com/schahab04-cpu/chula_class_ollama-codex_local_LLM} \\
\end{tabular}

\vspace{2em}

\section{Project Overview}

Yapperton is an oral study buddy that you can talk with. You upload your course
material, like lecture slides or PDFs, and then you can ask questions about it
and get easy summaries back, either in text or with your voice. Everything runs
locally on the own computer with Ollama, so no documents is sent away to some
company in the cloud. The idea came from that we all three study better when we
explain things out loud to each other, but it is not always that you have a
friend who has time to listen at 11 in the evening the day before the exam.

\subsection*{What problem does your project have?}

To read long course documents alone is honestly quite heavy, and many students
(us included) lose the focus after a few pages. It is easier to learn when you
can talk about the material and ask follow up questions, but a study partner is
not always available. The AI tools that already exist are mostly text based and
cloud based, which means you have to type a lot and also upload your material
to someone else's server. We wanted something where you can just speak and get
a short and simple explanation back, a bit like calling a classmate.

\subsection*{What are the main objectives of the project?}

\begin{itemize}[leftmargin=*]
  \item Build a chat where you can ask questions about your own uploaded
        documents (PDF, DOCX, TXT and MD) and get answers in plain language,
        with the sources shown so you can check what the answer builds on.
  \item Make a hands-free voice mode where you talk and the study buddy talks
        back, using Whisper for speech to text and Kokoro for text to speech.
  \item Keep everything local with a small model (qwen3:1.7b) so it works
        on a normal laptop and the material stays private.
  \item Make the summaries short and easy, so it feels more like a
        conversation than reading a textbook one more time.
\end{itemize}

\subsection*{What is the expected outcome?}

The expected outcome is a working web app that a student can start on their own
laptop, upload the course material in, and then have a spoken conversation with
about the content. We expect that it will make repetition before exams more fun
and less heavy, especially for people who learn better by listening and
talking. We are also aware that a small model can sometimes say wrong things,
so the goal is not to replace the lectures but to be a help on the side, and
that is why the sources always are shown in the text chat.

\end{document}
